\section{A local star obstruction for critical seven-pole profiles}
\label{sec:path-seven-star-obstruction}

The star-cap inequalities exclude one possible equality profile for a seven-cut.
This does not yet exclude every critical seven-pole profile.

\begin{lemma}[Two disjoint triangle covers on five vertices]
\label{lem:path-seven-triangle-covers}
Let $H$ be a graph on five vertices, with $d=10-|E(H)|$ missing edges.
For each edge $e$, let $x_e,y_e\in[0,1]$ satisfy $x_ey_e=0$.
Suppose every triangle $Q$ of $H$ satisfies
\[
 \sum_{e\in E(Q)}x_e\ge1,
 \qquad
 \sum_{e\in E(Q)}y_e\ge1.
\]
Then
\[
 \sum_{e\in E(H)}(x_e+y_e)\ge L(d),
 \qquad
 L(0)=10,\quad L(1)=6,\quad L(2)=4,\quad L(3)=3,
\]
where $L(d)=0$ for $d\ge4$.
In particular, the sum is at least $10-4d$.
\end{lemma}

\begin{proof}
First consider a triangle on vertices $u,v,w$.
The disjoint-support condition implies that one of its two covers has only one positive edge.
After exchanging the covers and relabeling the triangle, write
\[
 x_{uv}=1,\qquad x_{uw}=x_{vw}=y_{uv}=0,
 \qquad y_{uw}+y_{vw}\ge1.
\]
If another vertex $a$ is adjacent to all three triangle vertices, the triangles $auw$ and $avw$ require
\[
 x_{au}+x_{aw}\ge1,\qquad x_{av}+x_{aw}\ge1.
\]
If $x_{aw}<1$, both $x_{au},x_{av}$ are positive.
Then $y_{au}=y_{av}=0$, contradicting the $y$-cover of $auv$.
Hence $x_{aw}=1$.
The triangle $auv$ also gives $y_{au}+y_{av}\ge1$.
Thus the three edges from $a$ to the triangle have total mass at least two.

If $d=1$, take $a,b$ as the endpoints of the missing edge.
Their other three vertices form the triangle above.
Its mass is at least two, and each of $a,b$ contributes at least two.
This proves the lower bound six.

If $d=0$, the same argument gives $x_{aw}=x_{bw}=1$ and therefore $y_{ab}=1$.
The triangles $abu$ and $abv$ then require
\[
 x_{au}+x_{bu}\ge1,\qquad x_{av}+x_{bv}\ge1.
\]
Neither $a$ nor $b$ can have positive $x$-mass toward both $u$ and $v$, because its $y$-mass toward them is at least one.
Consequently these positive $x$-edges form a perfect matching between $\{a,b\}$ and $\{u,v\}$, with both weights one.
Assume that this matching is $au,bv$.
Then $y_{av}=y_{bu}=1$.
The triangles $auw$ and $bvw$ force $y_{uw}=y_{vw}=1$.
Every edge now has total mass one, proving the lower bound ten.

If $d=2$, the two missing edges are adjacent or disjoint.
For missing edges $01,02$, use the edge-disjoint triangles $034$ and $123$.
For missing edges $01,23$, use the edge-disjoint triangles $024$ and $134$.
Each triangle has total mass at least two, proving the bound four.

If $d=3$, the graph has at least two triangles.
To see this, a triangle-free graph on five vertices has at most six edges.
Indeed, a vertex of degree at least three has independent neighbors, and the remaining edges give this bound directly.
If there were exactly one triangle, each of the other two vertices could have at most one neighbor on it.
The graph would again have at most six edges.

Choose two triangles.
If they are edge-disjoint, their combined mass is at least four.
Otherwise they share one edge, on which one cover has mass zero.
That cover needs mass at least two on their disjoint remaining edges.
The other cover needs mass at least one in their union.
Their combined mass is therefore at least three.
For $d\ge4$, nonnegativity completes the asserted bounds.
\end{proof}

\begin{proposition}[A critical seven-cut needs a six-terminal sector]
\label{prop:path-seven-critical-six-sector}
Consider two limiting normalized seven-pole forest profiles obtained from shore ranks tending to infinity inside cubic graphs without cyclic cuts of size at most four.
If their normalized crossing functional equals $32$, at least one six-terminal forest total is nonzero.
\end{proposition}

\begin{proof}
The fixed-port aggregate bound is $32$ on each shore, and every limiting path mass is at most $1/2$.
For a profile, let $S$ be its total path mass and let
\[
 E=\sum_e a_e^2+\sum_{|T|=4}U_T^2+\sum_{|T|=6}W_T^2.
\]
Every forest total is at most one, so $E\le32-S/2$.
The crossing functional is at most the inner product of the two sector-total vectors, hence at most their average energy.
Equality at $32$ forces all path masses to vanish and all forest totals to be zero or one.
It also forces the two shores to have the same active sectors, exactly thirty-two of them.
For each active four-sector, corresponding pairing masses on opposite shores have disjoint supports.
Indeed, equal four-terminal pairings are precisely the incompatible pairs in that sector's crossing formula.

Suppose every six-terminal total is zero.
Then exactly thirty-two of the thirty-five four-terminal sectors are active.
Call the other three sectors missing.
Fix a terminal pair $e$ and form a graph $H_e$ on the other five ports.
An edge $f$ belongs to $H_e$ when the four-sector $e\cup f$ is active.
Let $x_f,y_f$ be the two shores' pairing masses for $e|f$.
They have disjoint supports.

If $Q$ is a triangle of $H_e$, all three relevant four-sector totals equal one.
The star-cap inequality on $e$ and $V(Q)$ gives
\[
 \sum_{f\in E(Q)}x_f\ge1,
 \qquad
 \sum_{f\in E(Q)}y_f\ge1.
\]
Thus Lemma~\ref{lem:path-seven-triangle-covers} applies.
Write $d_e=10-|E(H_e)|$.
Each missing four-sector contributes to $d_e$ for its six terminal pairs, so
\[
 \sum_e d_e=18.
\]
Summing the local lower bound gives
\[
 \sum_e\sum_{f\in E(H_e)}(x_f+y_f)
 \ge21\cdot10-4\cdot18=138.
\]
However, each pairing coordinate appears twice in this sum, once for each of its pairs.
Each active four-sector has total one on each shore.
Therefore the same sum equals $4\cdot32=128$, a contradiction.
\end{proof}

The stronger proposed local bound $10-2d$ is false.
For $K_5$ with one missing edge, two disjoint triangle covers can have total mass six.
The exploratory script \texttt{path-six-seven-local-star.py} records such a zero-one witness and checks all thirty-four graph types on five vertices numerically.
The proof above does not depend on that numerical search.

\subsection{One active six-terminal sector is also impossible}

\begin{lemma}[Four missing sectors with a sparse vertex]
\label{lem:path-seven-four-missing-cost}
Let $A_1,A_2,A_3,A_4$ be distinct four-subsets of a seven-element set.
Suppose some element $v$ belongs to at most one of them.
For each pair $e$, let $d_e$ count the sets $A_i$ containing $e$.
Then
\[
 \sum_e L(d_e)\ge126,
\]
where $L$ is defined in Lemma~\ref{lem:path-seven-triangle-covers}.
\end{lemma}

\begin{proof}
Define
\[
 f(d)=L(d)-10+4d,
 \qquad (f(0),f(1),f(2),f(3),f(4))=(0,0,2,5,6).
\]
Since $\sum_e d_e=4\binom42=24$, it suffices to prove
\[
 \sum_e f(d_e)\ge12.
\]

Relabel the sets so that $A_1,A_2,A_3$ omit $v$.
Their complements in the other six elements form the three edges of a graph $H$.
For a pair $e=ab$ on those six elements, its codegree in $A_1,A_2,A_3$ is
\[
 q_e=3-\deg_H(a)-\deg_H(b)+\mathbf1_{ab\in E(H)}.
\]
The five possible types of $H$, with isolated vertices omitted from their names, give
\[
\begin{array}{c|ccccc}
 H&3K_2&P_3+K_2&P_4&K_{1,3}&K_3\\\hline
 \sum_e f(q_e)&6&10&13&17&15
\end{array}
\]
These values follow by substituting the degree sequence and the three edges into the displayed formula.

Put $D=A_4\setminus\{v\}$, so $|D|$ is three or four.
Adding $A_4$ changes $q_e$ to $q_e+1$ exactly for the pairs in $D$.
Pairs containing $v$ contribute zero to $f$, whether or not $A_4$ contains $v$.
The increments
\[
 f(q+1)-f(q)=0,2,3,1\qquad(q=0,1,2,3)
\]
are all nonnegative.

If $H=3K_2$, every increment is at least two.
Any three-element subset of $D$ therefore adds at least six to the base value six.
If $H=P_3+K_2$, exactly two pair increments are zero.
They join the center of $P_3$ to the two endpoints of $K_2$.
All other pair increments are at least two.
Thus every three-element subset of $D$ has total increment at least two, which raises the base value ten to at least twelve.
For the other three types, the base value already exceeds twelve.
This proves the lemma.
\end{proof}

\begin{corollary}[At least two active six-terminal sectors]
\label{cor:path-seven-critical-two-six-sectors}
Under the limiting-profile hypotheses of Proposition~\ref{prop:path-seven-critical-six-sector}, crossing equality at $32$ requires at least two active six-terminal sectors.
\end{corollary}

\begin{proof}
Proposition~\ref{prop:path-seven-critical-six-sector} excludes zero active six-terminal sectors.
Suppose exactly one such sector is active, on the six-set $T$.
There are then thirty-one active four-sectors and four missing four-sectors.
Lemma~\ref{lem:four-seven-active-six-three-missing} says that at least three of the missing four-sectors lie inside $T$.
Thus the omitted port belongs to at most one missing sector.

Apply Lemma~\ref{lem:path-seven-four-missing-cost} to the four missing sectors.
Summing the local triangle-cover bounds gives total pairing mass at least $126$.
The actual total is $4\cdot31=124$, by the same double-counting used above.
This contradiction excludes exactly one active six-terminal sector.
\end{proof}
